% Capítulo para incorporar a la tesis mediante % Capítulo para incorporar a la tesis mediante % Capítulo para incorporar a la tesis mediante % Capítulo para incorporar a la tesis mediante \input{capitulo_08}.
% Corte documental: 4 de noviembre de 2026.
% La capacitación y el despliegue institucional se presentan como planificación.
\chapter{Implantación y puesta en marcha}
\label{cap:implantacion-puesta-marcha}

El presente capítulo describe la preparación de los usuarios y las actividades
necesarias para incorporar F.A.M.A. al trabajo de un entorno de investigación.
La implantación comprende la configuración de los recursos tecnológicos, la
disponibilidad de los datos y modelos, la definición de responsabilidades y la
verificación de las operaciones que se realizarán mediante la plataforma.
Debido a su arquitectura híbrida, estas actividades deben coordinar el
almacenamiento centralizado en Google Cloud Platform con el procesamiento
ejecutado en una estación o servidor local.

En primer lugar, se establece el objetivo del plan de capacitación y se
organizan las actividades de aprendizaje para investigadores y administradores.
Posteriormente, se presenta una estrategia gradual de implantación, tomando
como antecedente la puesta en marcha del nodo local del proyecto. El plan
formativo se plantea como una propuesta de ejecución; la operación local
verificada y las condiciones pendientes para el despliegue institucional se
describen según su alcance.

\section{Objetivo del Plan de Capacitación}
\label{sec:objetivo-capacitacion}

El objetivo del plan de capacitación es desarrollar las competencias necesarias
para que los usuarios autorizados operen F.A.M.A. de manera autónoma y
responsable, comprendiendo el recorrido de los datos desde su incorporación
hasta la obtención y revisión de una predicción acústica. La formación busca
que los participantes utilicen la interfaz web, interpreten los estados de las
operaciones y conserven la trazabilidad de los experimentos.

Para alcanzar este objetivo, se plantean los siguientes resultados de
aprendizaje:

\begin{enumerate}
    \item Reconocer la distribución de responsabilidades entre el nodo de
    procesamiento local, PostgreSQL y Google Cloud Storage.
    \item Acceder a la plataforma con una cuenta propia y utilizar los recursos
    disponibles de acuerdo con el rol y las asignaciones del usuario.
    \item Identificar un conjunto de datos, comprobar su disponibilidad local y
    preparar las particiones que utilizará un experimento.
    \item Configurar y supervisar un entrenamiento individual, distinguir su
    inicio de su finalización y consultar los resultados registrados.
    \item Realizar predicciones, interpretar sus resultados y registrar
    validaciones o correcciones con su correspondiente procedencia.
    \item Reconocer incidentes habituales y aplicar las acciones de consulta,
    respaldo o derivación que correspondan a cada responsabilidad.
\end{enumerate}

La capacitación considera dos perfiles de uso de la aplicación. El
\textbf{investigador} se concentra en los conjuntos de datos asignados, los
experimentos y la revisión de predicciones. El \textbf{administrador} incorpora
la gestión de cuentas y asignaciones, la revisión de solicitudes y la
activación explícita de modelos, además de las operaciones investigativas que
tenga disponibles.

La preparación de la infraestructura corresponde al \textbf{responsable
técnico del nodo}. Esta función incluye la configuración de Google Cloud, las
credenciales técnicas, las migraciones y la continuidad del servicio. Es una
responsabilidad operacional que puede asumir una persona del equipo o de la
institución, y no constituye un tercer rol web. Los investigadores que acceden
únicamente desde el navegador no requieren recibir credenciales de Google
Cloud ni contraseñas de PostgreSQL.

\section{Fases del Plan de Entrenamiento}
\label{sec:fases-entrenamiento-usuarios}

En esta sección, el término entrenamiento se refiere a la formación práctica de
los usuarios de la plataforma. El entrenamiento de redes neuronales se aborda
como una de las operaciones que los participantes deben aprender a ejecutar y
supervisar.

Se propone una modalidad guiada que combine demostraciones, ejercicios sobre
datos controlados y revisión de resultados. El equipo de desarrollo estará a
cargo de la explicación inicial y del acompañamiento, mientras que el
responsable técnico preparará las cuentas, los permisos y el entorno de
práctica. Los materiales de apoyo comprenderán el manual de uso, las guías de
instalación y recuperación, y un conjunto de audios con etiquetas conocidas.

La Tabla~\ref{tab:plan-capacitacion} presenta una distribución inicial de cinco
horas. Las duraciones son estimaciones de planificación, sujetas a ajuste según
la experiencia de los participantes. El tiempo de cómputo de un entrenamiento
puede exceder el bloque formativo y deberá programarse considerando el equipo
disponible.

\begin{table}[htbp]
    \centering
    \caption{Distribución propuesta del plan de capacitación.}
    \label{tab:plan-capacitacion}
    \begingroup
    \small
    \renewcommand{\arraystretch}{1.18}
    \begin{tabular}{p{0.08\linewidth} p{0.39\linewidth} p{0.27\linewidth} p{0.10\linewidth}}
        \hline
        \textbf{Fase} & \textbf{Actividad} & \textbf{Participantes} & \textbf{Minutos} \\
        \hline
        1 & Inducción y reconocimiento de la plataforma & Investigadores y administradores & 30 \\
        2 & Acceso y gestión de datos & Investigadores y administradores & 60 \\
        3 & Entrenamiento y seguimiento de modelos & Investigadores y administradores & 75 \\
        4 & Inferencia y retroalimentación & Investigadores y administradores & 60 \\
        5 & Administración y continuidad operacional & Administradores y responsable técnico & 45 \\
        6 & Ejercicio integrado y evaluación & Participantes, según su responsabilidad & 30 \\
        \hline
        \multicolumn{3}{r}{\textbf{Total propuesto}} & \textbf{300} \\
        \hline
    \end{tabular}
    \endgroup
\end{table}

\subsection{Fase 1: Inducción y reconocimiento de la plataforma}

Se presentará el propósito de F.A.M.A. y su arquitectura híbrida. Los
participantes reconocerán las vistas de gestión de audios, entrenamiento,
predicción y consulta de resultados, junto con las funciones administrativas
que correspondan a su cuenta.

La demostración explicará la diferencia entre un archivo de audio, sus
metadatos y los pesos de un modelo. También se abordará la separación entre
almacenamiento cloud y cómputo local. Como resultado, cada participante deberá
identificar dónde se conservan los datos y qué equipo ejecuta el
procesamiento.

\subsection{Fase 2: Acceso y gestión de datos}

Se practicará el inicio y cierre de sesión, la consulta de recursos asignados y
la identificación del conjunto de datos utilizado en el ejercicio. Se
revisarán las operaciones de carga o sincronización disponibles, sus estados y
las diferencias entre archivos registrados, archivos presentes en GCS y
archivos descargados en el nodo.

El ejercicio incluirá la revisión de clases y metadatos, seguida de la
previsualización de particiones de entrenamiento, validación y prueba. Se
explicará por qué grabaciones relacionadas deben conservarse dentro de una
misma partición y por qué añadir audios exige preparar una versión nueva. El
resultado esperado será reconocer una fuente local utilizable y una
preparación identificada para el experimento.

\subsection{Fase 3: Entrenamiento y seguimiento de modelos}

Se configurará un entrenamiento individual con un conjunto reducido y
preparado previamente. La actividad explicará el efecto de la arquitectura,
las épocas, el tamaño de lote y la tasa de aprendizaje, junto con las
restricciones impuestas por la memoria y el dispositivo de cómputo.

Los participantes consultarán el progreso, los registros y las métricas, y
distinguirán entre una solicitud admitida, un trabajo en ejecución y un
resultado final. Se revisará la detención cooperativa cuando sea aplicable,
sin interpretarla como una cancelación instantánea ni como una garantía de
guardar el componente en curso.

El resultado esperado será localizar el modelo registrado y sus evaluaciones
una vez concluido el trabajo. Se explicará que completar una evaluación
habilita la revisión del artefacto, mientras que su activación como
predeterminado requiere una decisión administrativa explícita.

\subsection{Fase 4: Inferencia y retroalimentación}

Se utilizará un modelo disponible para clasificar audios de prueba con
etiquetas conocidas. El participante revisará el archivo seleccionado, el
modelo utilizado y la respuesta obtenida. La interpretación incluirá las
salidas de silencio o ruido y la distinción entre confianza reportada y
corrección efectiva de la clasificación.

Posteriormente, se registrará una validación o corrección y se revisará su
estado de incorporación. Se explicará que aprobar un aporte no modifica los
pesos de un modelo existente: su utilización requiere una preparación posterior
que lo incluya y un nuevo entrenamiento. La práctica se considerará completada
cuando el participante pueda relacionar el audio, la predicción y la
retroalimentación registrada.

\subsection{Fase 5: Administración y continuidad operacional}

El módulo administrativo abordará la gestión de cuentas, la asignación de
datasets y la selección del modelo predeterminado. El responsable técnico
demostrará cómo consultar la disponibilidad del nodo, reconocer fallos de
conexión y localizar los registros operativos.

Se revisarán los procedimientos de respaldo de checkpoints y de recuperación
de la base de datos, conservando la separación entre credenciales web,
identidad técnica cloud y usuarios SQL. Las restauraciones de práctica se
realizarán en destinos de prueba preparados para ese propósito. El resultado
esperado será que cada responsable identifique qué acciones puede realizar y
cuáles debe derivar.

\subsection{Fase 6: Ejercicio integrado y evaluación}

La evaluación será práctica y utilizará una lista de verificación acorde con
el perfil del participante. Se comprobará el acceso con cuenta propia, la
identificación del dataset y del experimento, la consulta de resultados y el
registro de una predicción con su revisión. Para los responsables de
administración se añadirá una tarea de gestión de recursos o de diagnóstico.

Se registrarán las tareas completadas, las dificultades y la asistencia
requerida. Las operaciones que permanezcan en ejecución se revisarán al
finalizar su cómputo. Las dificultades que afecten permisos, procedencia de
datos o interpretación de estados deberán resolverse antes de autorizar el
uso autónomo del recorrido correspondiente.

La impartición de las sesiones, la asistencia y sus resultados deberán
documentarse al ejecutar el plan. Las actividades y duraciones propuestas no
representan una capacitación ya realizada.

\section{Estrategia de Implantación y Puesta en Marcha}
\label{sec:estrategia-implantacion}

La estrategia considera una implantación gradual: preparación de
infraestructura, habilitación de un nodo controlado, verificación operacional
y posterior transferencia al entorno institucional. Esta secuencia permite
revisar las dependencias de la arquitectura híbrida y preservar los datos
existentes antes de ampliar el acceso.

El nodo local del proyecto constituye el antecedente de puesta en marcha. El
paso hacia un servidor de la institución requiere acordar su administración,
conectividad, continuidad y respaldo. La disponibilidad comprobada en el
equipo del proyecto no acredita por sí sola esas condiciones en otro destino.

\subsection{Preparación de la infraestructura y las responsabilidades}

La institución deberá designar el equipo que ejecutará el backend y la
interfaz web, además del responsable de su operación. Los investigadores
podrán utilizar navegadores como clientes; el procesamiento se realizará en
el servidor o estación configurada. La ejecución en GPU dependerá de la
disponibilidad y compatibilidad de CUDA y PyTorch; el recorrido implementado
también permite ejecución en CPU.

En Google Cloud se deberán identificar el proyecto, la instancia Cloud SQL,
la base PostgreSQL y el bucket de GCS. Se preparará una identidad técnica
autorizada para los recursos requeridos y se conservarán separados el usuario
SQL diario y el usuario que aplica migraciones. La conexión PostgreSQL del
entorno oficial utiliza Cloud SQL Auth Proxy; el acceso a objetos de GCS se
realiza mediante las bibliotecas de Google.

Antes del arranque se verificarán las dependencias, el espacio local y la
conectividad. Cada entorno conservará su perfil privado de instalación y una
raíz persistente propia para audios, particiones, checkpoints y registros. Las
credenciales no formarán parte del repositorio ni se distribuirán a los
investigadores que utilizan exclusivamente el navegador.

\subsection{Instalación, configuración y preparación de datos}

La instalación deberá utilizar una versión identificada del software y un
esquema SQL compatible. Las migraciones se ejecutarán como una operación
explícita del responsable autorizado; el servicio diario comprobará el estado
de la base sin asumir permisos para modificar su estructura.

Para una instalación nueva se utilizará el asistente protegido de
configuración, que conecta recursos previamente preparados y permite crear la
primera cuenta administradora. Si el destino contiene historia y una cuenta
administradora existente, se aplicará el procedimiento de adopción
correspondiente. Este tratamiento conserva las identidades previas y evita
repetir el alta inicial.

Los audios, índices y pesos se transferirán mediante una copia controlada o
recuperación desde GCS. Se comprobarán los tamaños y las huellas de los
artefactos antes de utilizarlos. Los modelos ya entrenados podrán incorporarse
al nodo sin repetir su entrenamiento, siempre que se conserven su
compatibilidad, configuración y clases.

La distribución del código no sustituye esta preparación de archivos. Del
mismo modo, la existencia de un registro en PostgreSQL no demuestra que el
checkpoint requerido esté presente en el nuevo equipo.

\subsection{Arranque y verificación operacional}

El arranque comprenderá la conexión cloud, el backend y la interfaz web con
su comunicación configurada. Se comprobarán por separado la respuesta del
servicio, la compatibilidad del esquema y el acceso a los recursos. Después se
verificarán el inicio y cierre de sesión, los permisos, el catálogo y la
disponibilidad local de los datos y pesos seleccionados.

La aceptación del recorrido funcional deberá incluir:

\begin{enumerate}
    \item Acceso con cuentas propias y comprobación de los recursos
    autorizados para cada perfil.
    \item Consulta de un dataset y verificación de sus archivos locales,
    metadatos y particiones.
    \item Entrenamiento individual controlado, comprobación del checkpoint y
    consulta de las evaluaciones registradas.
    \item Inferencia con un modelo disponible, registro del resultado y
    seguimiento de una validación o corrección.
    \item Reinicio del nodo, conservación de su configuración y recuperación
    comprobada de los respaldos en un destino separado.
\end{enumerate}

Las pruebas deberán registrar la versión del código, los artefactos utilizados,
los resultados y las incidencias. Se emplearán datos de prueba o una selección
autorizada, evitando alterar conjuntos de evaluación reservados. El éxito de
un recorrido funcional no sustituye la evaluación científica de la calidad
del clasificador ni la medición de los requisitos de rendimiento.

\subsection{Puesta en marcha local documentada}

Con corte al 4 de noviembre de 2026, el proyecto dispone de un nodo local
configurado con PostgreSQL en Cloud SQL, almacenamiento en GCS y archivos
persistentes propios. La actualización del 25 de octubre aplicó la revisión
SQL \texttt{0004\_integration}, conservó la historia previa y comprobó el
acceso, los permisos y el catálogo de modelos. También se restauraron copias
lógicas anterior y posterior en bases locales separadas, comparando tablas
y secuencias.

El catálogo verificado en ese mantenimiento incluyó un ensamble de aves como
predeterminado y cinco modelos individuales. La preparación no acreditó la
disponibilidad del ensamble industrial, cuyos pesos seguían pendientes. La
operación local de aves constituye un punto de partida para la implantación y
no permite extender sus resultados a otros dominios sin comprobaciones propias.

El 1 de noviembre se documentó la reanudación de los servicios y la
comprobación del acceso administrador. Ese registro no estableció la causa de
su detención previa ni verificó un mecanismo permanente de arranque. El 4 de
noviembre se comprobó la cuenta y el permiso técnico de un segundo integrante
del equipo; el arranque de su nodo Linux permaneció pendiente.

La Tabla~\ref{tab:estado-implantacion} resume el alcance de estos antecedentes.

\begin{table}[htbp]
    \centering
    \caption{Estado documentado de la implantación al 4 de noviembre de 2026.}
    \label{tab:estado-implantacion}
    \begingroup
    \small
    \renewcommand{\arraystretch}{1.18}
    \begin{tabular}{p{0.22\linewidth} p{0.68\linewidth}}
        \hline
        \textbf{Aspecto} & \textbf{Alcance observado o pendiente} \\
        \hline
        Nodo local & Configuración oficial, conexiones cloud, acceso y catálogo comprobados en el equipo del proyecto. \\
        Recuperación SQL & Respaldos lógicos restaurados y contrastados en bases locales separadas. Política externa y recuperación administrada pendientes. \\
        Segundo equipo & Cuenta web y permiso técnico comprobados; configuración y arranque Linux pendientes de verificación. \\
        Capacitación & Plan formativo propuesto; ejecución y evaluación con participantes pendientes. \\
        Entorno institucional & Servidor, acceso HTTPS, servicio permanente y política de continuidad pendientes de preparación y aceptación. \\
        \hline
    \end{tabular}
    \endgroup
\end{table}

\subsection{Continuidad operacional y transferencia institucional}

La transferencia a la institución deberá formalizar quién administra las
cuentas, quién resuelve incidentes y quién aplica actualizaciones. Se
prepararán el origen web autorizado, HTTPS y la supervisión de los procesos,
incluido su comportamiento después de reiniciar el equipo.

El respaldo deberá considerar la base de datos y los archivos operativos
necesarios para recuperar el nodo. Se definirán destino, periodicidad,
retención y responsable de las copias, y se ensayará la restauración en un
entorno separado. El respaldo de un checkpoint no reemplaza la recuperación
de los registros SQL ni de los índices de datos.

Antes de una actualización se conservarán una copia verificada de la base y
la identificación del software y de los artefactos previos. Una reversión
deberá utilizar código y esquema compatibles; cambiar únicamente la versión
del repositorio no revierte una migración SQL.

La aceptación institucional se realizará cuando el entorno de destino
complete el recorrido operacional, el personal designado demuestre las
competencias correspondientes y se disponga de procedimientos de soporte y
recuperación comprobados. La importación y exportación ZIP se habilitarán
únicamente donde estén disponibles Linux/Docker y los workers requeridos.
Los resultados de esta aceptación deberán incorporarse al registro de puesta
en marcha cuando se ejecute la transferencia.
.
% Corte documental: 4 de noviembre de 2026.
% La capacitación y el despliegue institucional se presentan como planificación.
\chapter{Implantación y puesta en marcha}
\label{cap:implantacion-puesta-marcha}

El presente capítulo describe la preparación de los usuarios y las actividades
necesarias para incorporar F.A.M.A. al trabajo de un entorno de investigación.
La implantación comprende la configuración de los recursos tecnológicos, la
disponibilidad de los datos y modelos, la definición de responsabilidades y la
verificación de las operaciones que se realizarán mediante la plataforma.
Debido a su arquitectura híbrida, estas actividades deben coordinar el
almacenamiento centralizado en Google Cloud Platform con el procesamiento
ejecutado en una estación o servidor local.

En primer lugar, se establece el objetivo del plan de capacitación y se
organizan las actividades de aprendizaje para investigadores y administradores.
Posteriormente, se presenta una estrategia gradual de implantación, tomando
como antecedente la puesta en marcha del nodo local del proyecto. El plan
formativo se plantea como una propuesta de ejecución; la operación local
verificada y las condiciones pendientes para el despliegue institucional se
describen según su alcance.

\section{Objetivo del Plan de Capacitación}
\label{sec:objetivo-capacitacion}

El objetivo del plan de capacitación es desarrollar las competencias necesarias
para que los usuarios autorizados operen F.A.M.A. de manera autónoma y
responsable, comprendiendo el recorrido de los datos desde su incorporación
hasta la obtención y revisión de una predicción acústica. La formación busca
que los participantes utilicen la interfaz web, interpreten los estados de las
operaciones y conserven la trazabilidad de los experimentos.

Para alcanzar este objetivo, se plantean los siguientes resultados de
aprendizaje:

\begin{enumerate}
    \item Reconocer la distribución de responsabilidades entre el nodo de
    procesamiento local, PostgreSQL y Google Cloud Storage.
    \item Acceder a la plataforma con una cuenta propia y utilizar los recursos
    disponibles de acuerdo con el rol y las asignaciones del usuario.
    \item Identificar un conjunto de datos, comprobar su disponibilidad local y
    preparar las particiones que utilizará un experimento.
    \item Configurar y supervisar un entrenamiento individual, distinguir su
    inicio de su finalización y consultar los resultados registrados.
    \item Realizar predicciones, interpretar sus resultados y registrar
    validaciones o correcciones con su correspondiente procedencia.
    \item Reconocer incidentes habituales y aplicar las acciones de consulta,
    respaldo o derivación que correspondan a cada responsabilidad.
\end{enumerate}

La capacitación considera dos perfiles de uso de la aplicación. El
\textbf{investigador} se concentra en los conjuntos de datos asignados, los
experimentos y la revisión de predicciones. El \textbf{administrador} incorpora
la gestión de cuentas y asignaciones, la revisión de solicitudes y la
activación explícita de modelos, además de las operaciones investigativas que
tenga disponibles.

La preparación de la infraestructura corresponde al \textbf{responsable
técnico del nodo}. Esta función incluye la configuración de Google Cloud, las
credenciales técnicas, las migraciones y la continuidad del servicio. Es una
responsabilidad operacional que puede asumir una persona del equipo o de la
institución, y no constituye un tercer rol web. Los investigadores que acceden
únicamente desde el navegador no requieren recibir credenciales de Google
Cloud ni contraseñas de PostgreSQL.

\section{Fases del Plan de Entrenamiento}
\label{sec:fases-entrenamiento-usuarios}

En esta sección, el término entrenamiento se refiere a la formación práctica de
los usuarios de la plataforma. El entrenamiento de redes neuronales se aborda
como una de las operaciones que los participantes deben aprender a ejecutar y
supervisar.

Se propone una modalidad guiada que combine demostraciones, ejercicios sobre
datos controlados y revisión de resultados. El equipo de desarrollo estará a
cargo de la explicación inicial y del acompañamiento, mientras que el
responsable técnico preparará las cuentas, los permisos y el entorno de
práctica. Los materiales de apoyo comprenderán el manual de uso, las guías de
instalación y recuperación, y un conjunto de audios con etiquetas conocidas.

La Tabla~\ref{tab:plan-capacitacion} presenta una distribución inicial de cinco
horas. Las duraciones son estimaciones de planificación, sujetas a ajuste según
la experiencia de los participantes. El tiempo de cómputo de un entrenamiento
puede exceder el bloque formativo y deberá programarse considerando el equipo
disponible.

\begin{table}[htbp]
    \centering
    \caption{Distribución propuesta del plan de capacitación.}
    \label{tab:plan-capacitacion}
    \begingroup
    \small
    \renewcommand{\arraystretch}{1.18}
    \begin{tabular}{p{0.08\linewidth} p{0.39\linewidth} p{0.27\linewidth} p{0.10\linewidth}}
        \hline
        \textbf{Fase} & \textbf{Actividad} & \textbf{Participantes} & \textbf{Minutos} \\
        \hline
        1 & Inducción y reconocimiento de la plataforma & Investigadores y administradores & 30 \\
        2 & Acceso y gestión de datos & Investigadores y administradores & 60 \\
        3 & Entrenamiento y seguimiento de modelos & Investigadores y administradores & 75 \\
        4 & Inferencia y retroalimentación & Investigadores y administradores & 60 \\
        5 & Administración y continuidad operacional & Administradores y responsable técnico & 45 \\
        6 & Ejercicio integrado y evaluación & Participantes, según su responsabilidad & 30 \\
        \hline
        \multicolumn{3}{r}{\textbf{Total propuesto}} & \textbf{300} \\
        \hline
    \end{tabular}
    \endgroup
\end{table}

\subsection{Fase 1: Inducción y reconocimiento de la plataforma}

Se presentará el propósito de F.A.M.A. y su arquitectura híbrida. Los
participantes reconocerán las vistas de gestión de audios, entrenamiento,
predicción y consulta de resultados, junto con las funciones administrativas
que correspondan a su cuenta.

La demostración explicará la diferencia entre un archivo de audio, sus
metadatos y los pesos de un modelo. También se abordará la separación entre
almacenamiento cloud y cómputo local. Como resultado, cada participante deberá
identificar dónde se conservan los datos y qué equipo ejecuta el
procesamiento.

\subsection{Fase 2: Acceso y gestión de datos}

Se practicará el inicio y cierre de sesión, la consulta de recursos asignados y
la identificación del conjunto de datos utilizado en el ejercicio. Se
revisarán las operaciones de carga o sincronización disponibles, sus estados y
las diferencias entre archivos registrados, archivos presentes en GCS y
archivos descargados en el nodo.

El ejercicio incluirá la revisión de clases y metadatos, seguida de la
previsualización de particiones de entrenamiento, validación y prueba. Se
explicará por qué grabaciones relacionadas deben conservarse dentro de una
misma partición y por qué añadir audios exige preparar una versión nueva. El
resultado esperado será reconocer una fuente local utilizable y una
preparación identificada para el experimento.

\subsection{Fase 3: Entrenamiento y seguimiento de modelos}

Se configurará un entrenamiento individual con un conjunto reducido y
preparado previamente. La actividad explicará el efecto de la arquitectura,
las épocas, el tamaño de lote y la tasa de aprendizaje, junto con las
restricciones impuestas por la memoria y el dispositivo de cómputo.

Los participantes consultarán el progreso, los registros y las métricas, y
distinguirán entre una solicitud admitida, un trabajo en ejecución y un
resultado final. Se revisará la detención cooperativa cuando sea aplicable,
sin interpretarla como una cancelación instantánea ni como una garantía de
guardar el componente en curso.

El resultado esperado será localizar el modelo registrado y sus evaluaciones
una vez concluido el trabajo. Se explicará que completar una evaluación
habilita la revisión del artefacto, mientras que su activación como
predeterminado requiere una decisión administrativa explícita.

\subsection{Fase 4: Inferencia y retroalimentación}

Se utilizará un modelo disponible para clasificar audios de prueba con
etiquetas conocidas. El participante revisará el archivo seleccionado, el
modelo utilizado y la respuesta obtenida. La interpretación incluirá las
salidas de silencio o ruido y la distinción entre confianza reportada y
corrección efectiva de la clasificación.

Posteriormente, se registrará una validación o corrección y se revisará su
estado de incorporación. Se explicará que aprobar un aporte no modifica los
pesos de un modelo existente: su utilización requiere una preparación posterior
que lo incluya y un nuevo entrenamiento. La práctica se considerará completada
cuando el participante pueda relacionar el audio, la predicción y la
retroalimentación registrada.

\subsection{Fase 5: Administración y continuidad operacional}

El módulo administrativo abordará la gestión de cuentas, la asignación de
datasets y la selección del modelo predeterminado. El responsable técnico
demostrará cómo consultar la disponibilidad del nodo, reconocer fallos de
conexión y localizar los registros operativos.

Se revisarán los procedimientos de respaldo de checkpoints y de recuperación
de la base de datos, conservando la separación entre credenciales web,
identidad técnica cloud y usuarios SQL. Las restauraciones de práctica se
realizarán en destinos de prueba preparados para ese propósito. El resultado
esperado será que cada responsable identifique qué acciones puede realizar y
cuáles debe derivar.

\subsection{Fase 6: Ejercicio integrado y evaluación}

La evaluación será práctica y utilizará una lista de verificación acorde con
el perfil del participante. Se comprobará el acceso con cuenta propia, la
identificación del dataset y del experimento, la consulta de resultados y el
registro de una predicción con su revisión. Para los responsables de
administración se añadirá una tarea de gestión de recursos o de diagnóstico.

Se registrarán las tareas completadas, las dificultades y la asistencia
requerida. Las operaciones que permanezcan en ejecución se revisarán al
finalizar su cómputo. Las dificultades que afecten permisos, procedencia de
datos o interpretación de estados deberán resolverse antes de autorizar el
uso autónomo del recorrido correspondiente.

La impartición de las sesiones, la asistencia y sus resultados deberán
documentarse al ejecutar el plan. Las actividades y duraciones propuestas no
representan una capacitación ya realizada.

\section{Estrategia de Implantación y Puesta en Marcha}
\label{sec:estrategia-implantacion}

La estrategia considera una implantación gradual: preparación de
infraestructura, habilitación de un nodo controlado, verificación operacional
y posterior transferencia al entorno institucional. Esta secuencia permite
revisar las dependencias de la arquitectura híbrida y preservar los datos
existentes antes de ampliar el acceso.

El nodo local del proyecto constituye el antecedente de puesta en marcha. El
paso hacia un servidor de la institución requiere acordar su administración,
conectividad, continuidad y respaldo. La disponibilidad comprobada en el
equipo del proyecto no acredita por sí sola esas condiciones en otro destino.

\subsection{Preparación de la infraestructura y las responsabilidades}

La institución deberá designar el equipo que ejecutará el backend y la
interfaz web, además del responsable de su operación. Los investigadores
podrán utilizar navegadores como clientes; el procesamiento se realizará en
el servidor o estación configurada. La ejecución en GPU dependerá de la
disponibilidad y compatibilidad de CUDA y PyTorch; el recorrido implementado
también permite ejecución en CPU.

En Google Cloud se deberán identificar el proyecto, la instancia Cloud SQL,
la base PostgreSQL y el bucket de GCS. Se preparará una identidad técnica
autorizada para los recursos requeridos y se conservarán separados el usuario
SQL diario y el usuario que aplica migraciones. La conexión PostgreSQL del
entorno oficial utiliza Cloud SQL Auth Proxy; el acceso a objetos de GCS se
realiza mediante las bibliotecas de Google.

Antes del arranque se verificarán las dependencias, el espacio local y la
conectividad. Cada entorno conservará su perfil privado de instalación y una
raíz persistente propia para audios, particiones, checkpoints y registros. Las
credenciales no formarán parte del repositorio ni se distribuirán a los
investigadores que utilizan exclusivamente el navegador.

\subsection{Instalación, configuración y preparación de datos}

La instalación deberá utilizar una versión identificada del software y un
esquema SQL compatible. Las migraciones se ejecutarán como una operación
explícita del responsable autorizado; el servicio diario comprobará el estado
de la base sin asumir permisos para modificar su estructura.

Para una instalación nueva se utilizará el asistente protegido de
configuración, que conecta recursos previamente preparados y permite crear la
primera cuenta administradora. Si el destino contiene historia y una cuenta
administradora existente, se aplicará el procedimiento de adopción
correspondiente. Este tratamiento conserva las identidades previas y evita
repetir el alta inicial.

Los audios, índices y pesos se transferirán mediante una copia controlada o
recuperación desde GCS. Se comprobarán los tamaños y las huellas de los
artefactos antes de utilizarlos. Los modelos ya entrenados podrán incorporarse
al nodo sin repetir su entrenamiento, siempre que se conserven su
compatibilidad, configuración y clases.

La distribución del código no sustituye esta preparación de archivos. Del
mismo modo, la existencia de un registro en PostgreSQL no demuestra que el
checkpoint requerido esté presente en el nuevo equipo.

\subsection{Arranque y verificación operacional}

El arranque comprenderá la conexión cloud, el backend y la interfaz web con
su comunicación configurada. Se comprobarán por separado la respuesta del
servicio, la compatibilidad del esquema y el acceso a los recursos. Después se
verificarán el inicio y cierre de sesión, los permisos, el catálogo y la
disponibilidad local de los datos y pesos seleccionados.

La aceptación del recorrido funcional deberá incluir:

\begin{enumerate}
    \item Acceso con cuentas propias y comprobación de los recursos
    autorizados para cada perfil.
    \item Consulta de un dataset y verificación de sus archivos locales,
    metadatos y particiones.
    \item Entrenamiento individual controlado, comprobación del checkpoint y
    consulta de las evaluaciones registradas.
    \item Inferencia con un modelo disponible, registro del resultado y
    seguimiento de una validación o corrección.
    \item Reinicio del nodo, conservación de su configuración y recuperación
    comprobada de los respaldos en un destino separado.
\end{enumerate}

Las pruebas deberán registrar la versión del código, los artefactos utilizados,
los resultados y las incidencias. Se emplearán datos de prueba o una selección
autorizada, evitando alterar conjuntos de evaluación reservados. El éxito de
un recorrido funcional no sustituye la evaluación científica de la calidad
del clasificador ni la medición de los requisitos de rendimiento.

\subsection{Puesta en marcha local documentada}

Con corte al 4 de noviembre de 2026, el proyecto dispone de un nodo local
configurado con PostgreSQL en Cloud SQL, almacenamiento en GCS y archivos
persistentes propios. La actualización del 25 de octubre aplicó la revisión
SQL \texttt{0004\_integration}, conservó la historia previa y comprobó el
acceso, los permisos y el catálogo de modelos. También se restauraron copias
lógicas anterior y posterior en bases locales separadas, comparando tablas
y secuencias.

El catálogo verificado en ese mantenimiento incluyó un ensamble de aves como
predeterminado y cinco modelos individuales. La preparación no acreditó la
disponibilidad del ensamble industrial, cuyos pesos seguían pendientes. La
operación local de aves constituye un punto de partida para la implantación y
no permite extender sus resultados a otros dominios sin comprobaciones propias.

El 1 de noviembre se documentó la reanudación de los servicios y la
comprobación del acceso administrador. Ese registro no estableció la causa de
su detención previa ni verificó un mecanismo permanente de arranque. El 4 de
noviembre se comprobó la cuenta y el permiso técnico de un segundo integrante
del equipo; el arranque de su nodo Linux permaneció pendiente.

La Tabla~\ref{tab:estado-implantacion} resume el alcance de estos antecedentes.

\begin{table}[htbp]
    \centering
    \caption{Estado documentado de la implantación al 4 de noviembre de 2026.}
    \label{tab:estado-implantacion}
    \begingroup
    \small
    \renewcommand{\arraystretch}{1.18}
    \begin{tabular}{p{0.22\linewidth} p{0.68\linewidth}}
        \hline
        \textbf{Aspecto} & \textbf{Alcance observado o pendiente} \\
        \hline
        Nodo local & Configuración oficial, conexiones cloud, acceso y catálogo comprobados en el equipo del proyecto. \\
        Recuperación SQL & Respaldos lógicos restaurados y contrastados en bases locales separadas. Política externa y recuperación administrada pendientes. \\
        Segundo equipo & Cuenta web y permiso técnico comprobados; configuración y arranque Linux pendientes de verificación. \\
        Capacitación & Plan formativo propuesto; ejecución y evaluación con participantes pendientes. \\
        Entorno institucional & Servidor, acceso HTTPS, servicio permanente y política de continuidad pendientes de preparación y aceptación. \\
        \hline
    \end{tabular}
    \endgroup
\end{table}

\subsection{Continuidad operacional y transferencia institucional}

La transferencia a la institución deberá formalizar quién administra las
cuentas, quién resuelve incidentes y quién aplica actualizaciones. Se
prepararán el origen web autorizado, HTTPS y la supervisión de los procesos,
incluido su comportamiento después de reiniciar el equipo.

El respaldo deberá considerar la base de datos y los archivos operativos
necesarios para recuperar el nodo. Se definirán destino, periodicidad,
retención y responsable de las copias, y se ensayará la restauración en un
entorno separado. El respaldo de un checkpoint no reemplaza la recuperación
de los registros SQL ni de los índices de datos.

Antes de una actualización se conservarán una copia verificada de la base y
la identificación del software y de los artefactos previos. Una reversión
deberá utilizar código y esquema compatibles; cambiar únicamente la versión
del repositorio no revierte una migración SQL.

La aceptación institucional se realizará cuando el entorno de destino
complete el recorrido operacional, el personal designado demuestre las
competencias correspondientes y se disponga de procedimientos de soporte y
recuperación comprobados. La importación y exportación ZIP se habilitarán
únicamente donde estén disponibles Linux/Docker y los workers requeridos.
Los resultados de esta aceptación deberán incorporarse al registro de puesta
en marcha cuando se ejecute la transferencia.
.
% Corte documental: 4 de noviembre de 2026.
% La capacitación y el despliegue institucional se presentan como planificación.
\chapter{Implantación y puesta en marcha}
\label{cap:implantacion-puesta-marcha}

El presente capítulo describe la preparación de los usuarios y las actividades
necesarias para incorporar F.A.M.A. al trabajo de un entorno de investigación.
La implantación comprende la configuración de los recursos tecnológicos, la
disponibilidad de los datos y modelos, la definición de responsabilidades y la
verificación de las operaciones que se realizarán mediante la plataforma.
Debido a su arquitectura híbrida, estas actividades deben coordinar el
almacenamiento centralizado en Google Cloud Platform con el procesamiento
ejecutado en una estación o servidor local.

En primer lugar, se establece el objetivo del plan de capacitación y se
organizan las actividades de aprendizaje para investigadores y administradores.
Posteriormente, se presenta una estrategia gradual de implantación, tomando
como antecedente la puesta en marcha del nodo local del proyecto. El plan
formativo se plantea como una propuesta de ejecución; la operación local
verificada y las condiciones pendientes para el despliegue institucional se
describen según su alcance.

\section{Objetivo del Plan de Capacitación}
\label{sec:objetivo-capacitacion}

El objetivo del plan de capacitación es desarrollar las competencias necesarias
para que los usuarios autorizados operen F.A.M.A. de manera autónoma y
responsable, comprendiendo el recorrido de los datos desde su incorporación
hasta la obtención y revisión de una predicción acústica. La formación busca
que los participantes utilicen la interfaz web, interpreten los estados de las
operaciones y conserven la trazabilidad de los experimentos.

Para alcanzar este objetivo, se plantean los siguientes resultados de
aprendizaje:

\begin{enumerate}
    \item Reconocer la distribución de responsabilidades entre el nodo de
    procesamiento local, PostgreSQL y Google Cloud Storage.
    \item Acceder a la plataforma con una cuenta propia y utilizar los recursos
    disponibles de acuerdo con el rol y las asignaciones del usuario.
    \item Identificar un conjunto de datos, comprobar su disponibilidad local y
    preparar las particiones que utilizará un experimento.
    \item Configurar y supervisar un entrenamiento individual, distinguir su
    inicio de su finalización y consultar los resultados registrados.
    \item Realizar predicciones, interpretar sus resultados y registrar
    validaciones o correcciones con su correspondiente procedencia.
    \item Reconocer incidentes habituales y aplicar las acciones de consulta,
    respaldo o derivación que correspondan a cada responsabilidad.
\end{enumerate}

La capacitación considera dos perfiles de uso de la aplicación. El
\textbf{investigador} se concentra en los conjuntos de datos asignados, los
experimentos y la revisión de predicciones. El \textbf{administrador} incorpora
la gestión de cuentas y asignaciones, la revisión de solicitudes y la
activación explícita de modelos, además de las operaciones investigativas que
tenga disponibles.

La preparación de la infraestructura corresponde al \textbf{responsable
técnico del nodo}. Esta función incluye la configuración de Google Cloud, las
credenciales técnicas, las migraciones y la continuidad del servicio. Es una
responsabilidad operacional que puede asumir una persona del equipo o de la
institución, y no constituye un tercer rol web. Los investigadores que acceden
únicamente desde el navegador no requieren recibir credenciales de Google
Cloud ni contraseñas de PostgreSQL.

\section{Fases del Plan de Entrenamiento}
\label{sec:fases-entrenamiento-usuarios}

En esta sección, el término entrenamiento se refiere a la formación práctica de
los usuarios de la plataforma. El entrenamiento de redes neuronales se aborda
como una de las operaciones que los participantes deben aprender a ejecutar y
supervisar.

Se propone una modalidad guiada que combine demostraciones, ejercicios sobre
datos controlados y revisión de resultados. El equipo de desarrollo estará a
cargo de la explicación inicial y del acompañamiento, mientras que el
responsable técnico preparará las cuentas, los permisos y el entorno de
práctica. Los materiales de apoyo comprenderán el manual de uso, las guías de
instalación y recuperación, y un conjunto de audios con etiquetas conocidas.

La Tabla~\ref{tab:plan-capacitacion} presenta una distribución inicial de cinco
horas. Las duraciones son estimaciones de planificación, sujetas a ajuste según
la experiencia de los participantes. El tiempo de cómputo de un entrenamiento
puede exceder el bloque formativo y deberá programarse considerando el equipo
disponible.

\begin{table}[htbp]
    \centering
    \caption{Distribución propuesta del plan de capacitación.}
    \label{tab:plan-capacitacion}
    \begingroup
    \small
    \renewcommand{\arraystretch}{1.18}
    \begin{tabular}{p{0.08\linewidth} p{0.39\linewidth} p{0.27\linewidth} p{0.10\linewidth}}
        \hline
        \textbf{Fase} & \textbf{Actividad} & \textbf{Participantes} & \textbf{Minutos} \\
        \hline
        1 & Inducción y reconocimiento de la plataforma & Investigadores y administradores & 30 \\
        2 & Acceso y gestión de datos & Investigadores y administradores & 60 \\
        3 & Entrenamiento y seguimiento de modelos & Investigadores y administradores & 75 \\
        4 & Inferencia y retroalimentación & Investigadores y administradores & 60 \\
        5 & Administración y continuidad operacional & Administradores y responsable técnico & 45 \\
        6 & Ejercicio integrado y evaluación & Participantes, según su responsabilidad & 30 \\
        \hline
        \multicolumn{3}{r}{\textbf{Total propuesto}} & \textbf{300} \\
        \hline
    \end{tabular}
    \endgroup
\end{table}

\subsection{Fase 1: Inducción y reconocimiento de la plataforma}

Se presentará el propósito de F.A.M.A. y su arquitectura híbrida. Los
participantes reconocerán las vistas de gestión de audios, entrenamiento,
predicción y consulta de resultados, junto con las funciones administrativas
que correspondan a su cuenta.

La demostración explicará la diferencia entre un archivo de audio, sus
metadatos y los pesos de un modelo. También se abordará la separación entre
almacenamiento cloud y cómputo local. Como resultado, cada participante deberá
identificar dónde se conservan los datos y qué equipo ejecuta el
procesamiento.

\subsection{Fase 2: Acceso y gestión de datos}

Se practicará el inicio y cierre de sesión, la consulta de recursos asignados y
la identificación del conjunto de datos utilizado en el ejercicio. Se
revisarán las operaciones de carga o sincronización disponibles, sus estados y
las diferencias entre archivos registrados, archivos presentes en GCS y
archivos descargados en el nodo.

El ejercicio incluirá la revisión de clases y metadatos, seguida de la
previsualización de particiones de entrenamiento, validación y prueba. Se
explicará por qué grabaciones relacionadas deben conservarse dentro de una
misma partición y por qué añadir audios exige preparar una versión nueva. El
resultado esperado será reconocer una fuente local utilizable y una
preparación identificada para el experimento.

\subsection{Fase 3: Entrenamiento y seguimiento de modelos}

Se configurará un entrenamiento individual con un conjunto reducido y
preparado previamente. La actividad explicará el efecto de la arquitectura,
las épocas, el tamaño de lote y la tasa de aprendizaje, junto con las
restricciones impuestas por la memoria y el dispositivo de cómputo.

Los participantes consultarán el progreso, los registros y las métricas, y
distinguirán entre una solicitud admitida, un trabajo en ejecución y un
resultado final. Se revisará la detención cooperativa cuando sea aplicable,
sin interpretarla como una cancelación instantánea ni como una garantía de
guardar el componente en curso.

El resultado esperado será localizar el modelo registrado y sus evaluaciones
una vez concluido el trabajo. Se explicará que completar una evaluación
habilita la revisión del artefacto, mientras que su activación como
predeterminado requiere una decisión administrativa explícita.

\subsection{Fase 4: Inferencia y retroalimentación}

Se utilizará un modelo disponible para clasificar audios de prueba con
etiquetas conocidas. El participante revisará el archivo seleccionado, el
modelo utilizado y la respuesta obtenida. La interpretación incluirá las
salidas de silencio o ruido y la distinción entre confianza reportada y
corrección efectiva de la clasificación.

Posteriormente, se registrará una validación o corrección y se revisará su
estado de incorporación. Se explicará que aprobar un aporte no modifica los
pesos de un modelo existente: su utilización requiere una preparación posterior
que lo incluya y un nuevo entrenamiento. La práctica se considerará completada
cuando el participante pueda relacionar el audio, la predicción y la
retroalimentación registrada.

\subsection{Fase 5: Administración y continuidad operacional}

El módulo administrativo abordará la gestión de cuentas, la asignación de
datasets y la selección del modelo predeterminado. El responsable técnico
demostrará cómo consultar la disponibilidad del nodo, reconocer fallos de
conexión y localizar los registros operativos.

Se revisarán los procedimientos de respaldo de checkpoints y de recuperación
de la base de datos, conservando la separación entre credenciales web,
identidad técnica cloud y usuarios SQL. Las restauraciones de práctica se
realizarán en destinos de prueba preparados para ese propósito. El resultado
esperado será que cada responsable identifique qué acciones puede realizar y
cuáles debe derivar.

\subsection{Fase 6: Ejercicio integrado y evaluación}

La evaluación será práctica y utilizará una lista de verificación acorde con
el perfil del participante. Se comprobará el acceso con cuenta propia, la
identificación del dataset y del experimento, la consulta de resultados y el
registro de una predicción con su revisión. Para los responsables de
administración se añadirá una tarea de gestión de recursos o de diagnóstico.

Se registrarán las tareas completadas, las dificultades y la asistencia
requerida. Las operaciones que permanezcan en ejecución se revisarán al
finalizar su cómputo. Las dificultades que afecten permisos, procedencia de
datos o interpretación de estados deberán resolverse antes de autorizar el
uso autónomo del recorrido correspondiente.

La impartición de las sesiones, la asistencia y sus resultados deberán
documentarse al ejecutar el plan. Las actividades y duraciones propuestas no
representan una capacitación ya realizada.

\section{Estrategia de Implantación y Puesta en Marcha}
\label{sec:estrategia-implantacion}

La estrategia considera una implantación gradual: preparación de
infraestructura, habilitación de un nodo controlado, verificación operacional
y posterior transferencia al entorno institucional. Esta secuencia permite
revisar las dependencias de la arquitectura híbrida y preservar los datos
existentes antes de ampliar el acceso.

El nodo local del proyecto constituye el antecedente de puesta en marcha. El
paso hacia un servidor de la institución requiere acordar su administración,
conectividad, continuidad y respaldo. La disponibilidad comprobada en el
equipo del proyecto no acredita por sí sola esas condiciones en otro destino.

\subsection{Preparación de la infraestructura y las responsabilidades}

La institución deberá designar el equipo que ejecutará el backend y la
interfaz web, además del responsable de su operación. Los investigadores
podrán utilizar navegadores como clientes; el procesamiento se realizará en
el servidor o estación configurada. La ejecución en GPU dependerá de la
disponibilidad y compatibilidad de CUDA y PyTorch; el recorrido implementado
también permite ejecución en CPU.

En Google Cloud se deberán identificar el proyecto, la instancia Cloud SQL,
la base PostgreSQL y el bucket de GCS. Se preparará una identidad técnica
autorizada para los recursos requeridos y se conservarán separados el usuario
SQL diario y el usuario que aplica migraciones. La conexión PostgreSQL del
entorno oficial utiliza Cloud SQL Auth Proxy; el acceso a objetos de GCS se
realiza mediante las bibliotecas de Google.

Antes del arranque se verificarán las dependencias, el espacio local y la
conectividad. Cada entorno conservará su perfil privado de instalación y una
raíz persistente propia para audios, particiones, checkpoints y registros. Las
credenciales no formarán parte del repositorio ni se distribuirán a los
investigadores que utilizan exclusivamente el navegador.

\subsection{Instalación, configuración y preparación de datos}

La instalación deberá utilizar una versión identificada del software y un
esquema SQL compatible. Las migraciones se ejecutarán como una operación
explícita del responsable autorizado; el servicio diario comprobará el estado
de la base sin asumir permisos para modificar su estructura.

Para una instalación nueva se utilizará el asistente protegido de
configuración, que conecta recursos previamente preparados y permite crear la
primera cuenta administradora. Si el destino contiene historia y una cuenta
administradora existente, se aplicará el procedimiento de adopción
correspondiente. Este tratamiento conserva las identidades previas y evita
repetir el alta inicial.

Los audios, índices y pesos se transferirán mediante una copia controlada o
recuperación desde GCS. Se comprobarán los tamaños y las huellas de los
artefactos antes de utilizarlos. Los modelos ya entrenados podrán incorporarse
al nodo sin repetir su entrenamiento, siempre que se conserven su
compatibilidad, configuración y clases.

La distribución del código no sustituye esta preparación de archivos. Del
mismo modo, la existencia de un registro en PostgreSQL no demuestra que el
checkpoint requerido esté presente en el nuevo equipo.

\subsection{Arranque y verificación operacional}

El arranque comprenderá la conexión cloud, el backend y la interfaz web con
su comunicación configurada. Se comprobarán por separado la respuesta del
servicio, la compatibilidad del esquema y el acceso a los recursos. Después se
verificarán el inicio y cierre de sesión, los permisos, el catálogo y la
disponibilidad local de los datos y pesos seleccionados.

La aceptación del recorrido funcional deberá incluir:

\begin{enumerate}
    \item Acceso con cuentas propias y comprobación de los recursos
    autorizados para cada perfil.
    \item Consulta de un dataset y verificación de sus archivos locales,
    metadatos y particiones.
    \item Entrenamiento individual controlado, comprobación del checkpoint y
    consulta de las evaluaciones registradas.
    \item Inferencia con un modelo disponible, registro del resultado y
    seguimiento de una validación o corrección.
    \item Reinicio del nodo, conservación de su configuración y recuperación
    comprobada de los respaldos en un destino separado.
\end{enumerate}

Las pruebas deberán registrar la versión del código, los artefactos utilizados,
los resultados y las incidencias. Se emplearán datos de prueba o una selección
autorizada, evitando alterar conjuntos de evaluación reservados. El éxito de
un recorrido funcional no sustituye la evaluación científica de la calidad
del clasificador ni la medición de los requisitos de rendimiento.

\subsection{Puesta en marcha local documentada}

Con corte al 4 de noviembre de 2026, el proyecto dispone de un nodo local
configurado con PostgreSQL en Cloud SQL, almacenamiento en GCS y archivos
persistentes propios. La actualización del 25 de octubre aplicó la revisión
SQL \texttt{0004\_integration}, conservó la historia previa y comprobó el
acceso, los permisos y el catálogo de modelos. También se restauraron copias
lógicas anterior y posterior en bases locales separadas, comparando tablas
y secuencias.

El catálogo verificado en ese mantenimiento incluyó un ensamble de aves como
predeterminado y cinco modelos individuales. La preparación no acreditó la
disponibilidad del ensamble industrial, cuyos pesos seguían pendientes. La
operación local de aves constituye un punto de partida para la implantación y
no permite extender sus resultados a otros dominios sin comprobaciones propias.

El 1 de noviembre se documentó la reanudación de los servicios y la
comprobación del acceso administrador. Ese registro no estableció la causa de
su detención previa ni verificó un mecanismo permanente de arranque. El 4 de
noviembre se comprobó la cuenta y el permiso técnico de un segundo integrante
del equipo; el arranque de su nodo Linux permaneció pendiente.

La Tabla~\ref{tab:estado-implantacion} resume el alcance de estos antecedentes.

\begin{table}[htbp]
    \centering
    \caption{Estado documentado de la implantación al 4 de noviembre de 2026.}
    \label{tab:estado-implantacion}
    \begingroup
    \small
    \renewcommand{\arraystretch}{1.18}
    \begin{tabular}{p{0.22\linewidth} p{0.68\linewidth}}
        \hline
        \textbf{Aspecto} & \textbf{Alcance observado o pendiente} \\
        \hline
        Nodo local & Configuración oficial, conexiones cloud, acceso y catálogo comprobados en el equipo del proyecto. \\
        Recuperación SQL & Respaldos lógicos restaurados y contrastados en bases locales separadas. Política externa y recuperación administrada pendientes. \\
        Segundo equipo & Cuenta web y permiso técnico comprobados; configuración y arranque Linux pendientes de verificación. \\
        Capacitación & Plan formativo propuesto; ejecución y evaluación con participantes pendientes. \\
        Entorno institucional & Servidor, acceso HTTPS, servicio permanente y política de continuidad pendientes de preparación y aceptación. \\
        \hline
    \end{tabular}
    \endgroup
\end{table}

\subsection{Continuidad operacional y transferencia institucional}

La transferencia a la institución deberá formalizar quién administra las
cuentas, quién resuelve incidentes y quién aplica actualizaciones. Se
prepararán el origen web autorizado, HTTPS y la supervisión de los procesos,
incluido su comportamiento después de reiniciar el equipo.

El respaldo deberá considerar la base de datos y los archivos operativos
necesarios para recuperar el nodo. Se definirán destino, periodicidad,
retención y responsable de las copias, y se ensayará la restauración en un
entorno separado. El respaldo de un checkpoint no reemplaza la recuperación
de los registros SQL ni de los índices de datos.

Antes de una actualización se conservarán una copia verificada de la base y
la identificación del software y de los artefactos previos. Una reversión
deberá utilizar código y esquema compatibles; cambiar únicamente la versión
del repositorio no revierte una migración SQL.

La aceptación institucional se realizará cuando el entorno de destino
complete el recorrido operacional, el personal designado demuestre las
competencias correspondientes y se disponga de procedimientos de soporte y
recuperación comprobados. La importación y exportación ZIP se habilitarán
únicamente donde estén disponibles Linux/Docker y los workers requeridos.
Los resultados de esta aceptación deberán incorporarse al registro de puesta
en marcha cuando se ejecute la transferencia.
.
% Corte documental: 4 de noviembre de 2026.
% La capacitación y el despliegue institucional se presentan como planificación.
\chapter{Implantación y puesta en marcha}
\label{cap:implantacion-puesta-marcha}

El presente capítulo describe la preparación de los usuarios y las actividades
necesarias para incorporar F.A.M.A. al trabajo de un entorno de investigación.
La implantación comprende la configuración de los recursos tecnológicos, la
disponibilidad de los datos y modelos, la definición de responsabilidades y la
verificación de las operaciones que se realizarán mediante la plataforma.
Debido a su arquitectura híbrida, estas actividades deben coordinar el
almacenamiento centralizado en Google Cloud Platform con el procesamiento
ejecutado en una estación o servidor local.

En primer lugar, se establece el objetivo del plan de capacitación y se
organizan las actividades de aprendizaje para investigadores y administradores.
Posteriormente, se presenta una estrategia gradual de implantación, tomando
como antecedente la puesta en marcha del nodo local del proyecto. El plan
formativo se plantea como una propuesta de ejecución; la operación local
verificada y las condiciones pendientes para el despliegue institucional se
describen según su alcance.

\section{Objetivo del Plan de Capacitación}
\label{sec:objetivo-capacitacion}

El objetivo del plan de capacitación es desarrollar las competencias necesarias
para que los usuarios autorizados operen F.A.M.A. de manera autónoma y
responsable, comprendiendo el recorrido de los datos desde su incorporación
hasta la obtención y revisión de una predicción acústica. La formación busca
que los participantes utilicen la interfaz web, interpreten los estados de las
operaciones y conserven la trazabilidad de los experimentos.

Para alcanzar este objetivo, se plantean los siguientes resultados de
aprendizaje:

\begin{enumerate}
    \item Reconocer la distribución de responsabilidades entre el nodo de
    procesamiento local, PostgreSQL y Google Cloud Storage.
    \item Acceder a la plataforma con una cuenta propia y utilizar los recursos
    disponibles de acuerdo con el rol y las asignaciones del usuario.
    \item Identificar un conjunto de datos, comprobar su disponibilidad local y
    preparar las particiones que utilizará un experimento.
    \item Configurar y supervisar un entrenamiento individual, distinguir su
    inicio de su finalización y consultar los resultados registrados.
    \item Realizar predicciones, interpretar sus resultados y registrar
    validaciones o correcciones con su correspondiente procedencia.
    \item Reconocer incidentes habituales y aplicar las acciones de consulta,
    respaldo o derivación que correspondan a cada responsabilidad.
\end{enumerate}

La capacitación considera dos perfiles de uso de la aplicación. El
\textbf{investigador} se concentra en los conjuntos de datos asignados, los
experimentos y la revisión de predicciones. El \textbf{administrador} incorpora
la gestión de cuentas y asignaciones, la revisión de solicitudes y la
activación explícita de modelos, además de las operaciones investigativas que
tenga disponibles.

La preparación de la infraestructura corresponde al \textbf{responsable
técnico del nodo}. Esta función incluye la configuración de Google Cloud, las
credenciales técnicas, las migraciones y la continuidad del servicio. Es una
responsabilidad operacional que puede asumir una persona del equipo o de la
institución, y no constituye un tercer rol web. Los investigadores que acceden
únicamente desde el navegador no requieren recibir credenciales de Google
Cloud ni contraseñas de PostgreSQL.

\section{Fases del Plan de Entrenamiento}
\label{sec:fases-entrenamiento-usuarios}

En esta sección, el término entrenamiento se refiere a la formación práctica de
los usuarios de la plataforma. El entrenamiento de redes neuronales se aborda
como una de las operaciones que los participantes deben aprender a ejecutar y
supervisar.

Se propone una modalidad guiada que combine demostraciones, ejercicios sobre
datos controlados y revisión de resultados. El equipo de desarrollo estará a
cargo de la explicación inicial y del acompañamiento, mientras que el
responsable técnico preparará las cuentas, los permisos y el entorno de
práctica. Los materiales de apoyo comprenderán el manual de uso, las guías de
instalación y recuperación, y un conjunto de audios con etiquetas conocidas.

La Tabla~\ref{tab:plan-capacitacion} presenta una distribución inicial de cinco
horas. Las duraciones son estimaciones de planificación, sujetas a ajuste según
la experiencia de los participantes. El tiempo de cómputo de un entrenamiento
puede exceder el bloque formativo y deberá programarse considerando el equipo
disponible.

\begin{table}[htbp]
    \centering
    \caption{Distribución propuesta del plan de capacitación.}
    \label{tab:plan-capacitacion}
    \begingroup
    \small
    \renewcommand{\arraystretch}{1.18}
    \begin{tabular}{p{0.08\linewidth} p{0.39\linewidth} p{0.27\linewidth} p{0.10\linewidth}}
        \hline
        \textbf{Fase} & \textbf{Actividad} & \textbf{Participantes} & \textbf{Minutos} \\
        \hline
        1 & Inducción y reconocimiento de la plataforma & Investigadores y administradores & 30 \\
        2 & Acceso y gestión de datos & Investigadores y administradores & 60 \\
        3 & Entrenamiento y seguimiento de modelos & Investigadores y administradores & 75 \\
        4 & Inferencia y retroalimentación & Investigadores y administradores & 60 \\
        5 & Administración y continuidad operacional & Administradores y responsable técnico & 45 \\
        6 & Ejercicio integrado y evaluación & Participantes, según su responsabilidad & 30 \\
        \hline
        \multicolumn{3}{r}{\textbf{Total propuesto}} & \textbf{300} \\
        \hline
    \end{tabular}
    \endgroup
\end{table}

\subsection{Fase 1: Inducción y reconocimiento de la plataforma}

Se presentará el propósito de F.A.M.A. y su arquitectura híbrida. Los
participantes reconocerán las vistas de gestión de audios, entrenamiento,
predicción y consulta de resultados, junto con las funciones administrativas
que correspondan a su cuenta.

La demostración explicará la diferencia entre un archivo de audio, sus
metadatos y los pesos de un modelo. También se abordará la separación entre
almacenamiento cloud y cómputo local. Como resultado, cada participante deberá
identificar dónde se conservan los datos y qué equipo ejecuta el
procesamiento.

\subsection{Fase 2: Acceso y gestión de datos}

Se practicará el inicio y cierre de sesión, la consulta de recursos asignados y
la identificación del conjunto de datos utilizado en el ejercicio. Se
revisarán las operaciones de carga o sincronización disponibles, sus estados y
las diferencias entre archivos registrados, archivos presentes en GCS y
archivos descargados en el nodo.

El ejercicio incluirá la revisión de clases y metadatos, seguida de la
previsualización de particiones de entrenamiento, validación y prueba. Se
explicará por qué grabaciones relacionadas deben conservarse dentro de una
misma partición y por qué añadir audios exige preparar una versión nueva. El
resultado esperado será reconocer una fuente local utilizable y una
preparación identificada para el experimento.

\subsection{Fase 3: Entrenamiento y seguimiento de modelos}

Se configurará un entrenamiento individual con un conjunto reducido y
preparado previamente. La actividad explicará el efecto de la arquitectura,
las épocas, el tamaño de lote y la tasa de aprendizaje, junto con las
restricciones impuestas por la memoria y el dispositivo de cómputo.

Los participantes consultarán el progreso, los registros y las métricas, y
distinguirán entre una solicitud admitida, un trabajo en ejecución y un
resultado final. Se revisará la detención cooperativa cuando sea aplicable,
sin interpretarla como una cancelación instantánea ni como una garantía de
guardar el componente en curso.

El resultado esperado será localizar el modelo registrado y sus evaluaciones
una vez concluido el trabajo. Se explicará que completar una evaluación
habilita la revisión del artefacto, mientras que su activación como
predeterminado requiere una decisión administrativa explícita.

\subsection{Fase 4: Inferencia y retroalimentación}

Se utilizará un modelo disponible para clasificar audios de prueba con
etiquetas conocidas. El participante revisará el archivo seleccionado, el
modelo utilizado y la respuesta obtenida. La interpretación incluirá las
salidas de silencio o ruido y la distinción entre confianza reportada y
corrección efectiva de la clasificación.

Posteriormente, se registrará una validación o corrección y se revisará su
estado de incorporación. Se explicará que aprobar un aporte no modifica los
pesos de un modelo existente: su utilización requiere una preparación posterior
que lo incluya y un nuevo entrenamiento. La práctica se considerará completada
cuando el participante pueda relacionar el audio, la predicción y la
retroalimentación registrada.

\subsection{Fase 5: Administración y continuidad operacional}

El módulo administrativo abordará la gestión de cuentas, la asignación de
datasets y la selección del modelo predeterminado. El responsable técnico
demostrará cómo consultar la disponibilidad del nodo, reconocer fallos de
conexión y localizar los registros operativos.

Se revisarán los procedimientos de respaldo de checkpoints y de recuperación
de la base de datos, conservando la separación entre credenciales web,
identidad técnica cloud y usuarios SQL. Las restauraciones de práctica se
realizarán en destinos de prueba preparados para ese propósito. El resultado
esperado será que cada responsable identifique qué acciones puede realizar y
cuáles debe derivar.

\subsection{Fase 6: Ejercicio integrado y evaluación}

La evaluación será práctica y utilizará una lista de verificación acorde con
el perfil del participante. Se comprobará el acceso con cuenta propia, la
identificación del dataset y del experimento, la consulta de resultados y el
registro de una predicción con su revisión. Para los responsables de
administración se añadirá una tarea de gestión de recursos o de diagnóstico.

Se registrarán las tareas completadas, las dificultades y la asistencia
requerida. Las operaciones que permanezcan en ejecución se revisarán al
finalizar su cómputo. Las dificultades que afecten permisos, procedencia de
datos o interpretación de estados deberán resolverse antes de autorizar el
uso autónomo del recorrido correspondiente.

La impartición de las sesiones, la asistencia y sus resultados deberán
documentarse al ejecutar el plan. Las actividades y duraciones propuestas no
representan una capacitación ya realizada.

\section{Estrategia de Implantación y Puesta en Marcha}
\label{sec:estrategia-implantacion}

La estrategia considera una implantación gradual: preparación de
infraestructura, habilitación de un nodo controlado, verificación operacional
y posterior transferencia al entorno institucional. Esta secuencia permite
revisar las dependencias de la arquitectura híbrida y preservar los datos
existentes antes de ampliar el acceso.

El nodo local del proyecto constituye el antecedente de puesta en marcha. El
paso hacia un servidor de la institución requiere acordar su administración,
conectividad, continuidad y respaldo. La disponibilidad comprobada en el
equipo del proyecto no acredita por sí sola esas condiciones en otro destino.

\subsection{Preparación de la infraestructura y las responsabilidades}

La institución deberá designar el equipo que ejecutará el backend y la
interfaz web, además del responsable de su operación. Los investigadores
podrán utilizar navegadores como clientes; el procesamiento se realizará en
el servidor o estación configurada. La ejecución en GPU dependerá de la
disponibilidad y compatibilidad de CUDA y PyTorch; el recorrido implementado
también permite ejecución en CPU.

En Google Cloud se deberán identificar el proyecto, la instancia Cloud SQL,
la base PostgreSQL y el bucket de GCS. Se preparará una identidad técnica
autorizada para los recursos requeridos y se conservarán separados el usuario
SQL diario y el usuario que aplica migraciones. La conexión PostgreSQL del
entorno oficial utiliza Cloud SQL Auth Proxy; el acceso a objetos de GCS se
realiza mediante las bibliotecas de Google.

Antes del arranque se verificarán las dependencias, el espacio local y la
conectividad. Cada entorno conservará su perfil privado de instalación y una
raíz persistente propia para audios, particiones, checkpoints y registros. Las
credenciales no formarán parte del repositorio ni se distribuirán a los
investigadores que utilizan exclusivamente el navegador.

\subsection{Instalación, configuración y preparación de datos}

La instalación deberá utilizar una versión identificada del software y un
esquema SQL compatible. Las migraciones se ejecutarán como una operación
explícita del responsable autorizado; el servicio diario comprobará el estado
de la base sin asumir permisos para modificar su estructura.

Para una instalación nueva se utilizará el asistente protegido de
configuración, que conecta recursos previamente preparados y permite crear la
primera cuenta administradora. Si el destino contiene historia y una cuenta
administradora existente, se aplicará el procedimiento de adopción
correspondiente. Este tratamiento conserva las identidades previas y evita
repetir el alta inicial.

Los audios, índices y pesos se transferirán mediante una copia controlada o
recuperación desde GCS. Se comprobarán los tamaños y las huellas de los
artefactos antes de utilizarlos. Los modelos ya entrenados podrán incorporarse
al nodo sin repetir su entrenamiento, siempre que se conserven su
compatibilidad, configuración y clases.

La distribución del código no sustituye esta preparación de archivos. Del
mismo modo, la existencia de un registro en PostgreSQL no demuestra que el
checkpoint requerido esté presente en el nuevo equipo.

\subsection{Arranque y verificación operacional}

El arranque comprenderá la conexión cloud, el backend y la interfaz web con
su comunicación configurada. Se comprobarán por separado la respuesta del
servicio, la compatibilidad del esquema y el acceso a los recursos. Después se
verificarán el inicio y cierre de sesión, los permisos, el catálogo y la
disponibilidad local de los datos y pesos seleccionados.

La aceptación del recorrido funcional deberá incluir:

\begin{enumerate}
    \item Acceso con cuentas propias y comprobación de los recursos
    autorizados para cada perfil.
    \item Consulta de un dataset y verificación de sus archivos locales,
    metadatos y particiones.
    \item Entrenamiento individual controlado, comprobación del checkpoint y
    consulta de las evaluaciones registradas.
    \item Inferencia con un modelo disponible, registro del resultado y
    seguimiento de una validación o corrección.
    \item Reinicio del nodo, conservación de su configuración y recuperación
    comprobada de los respaldos en un destino separado.
\end{enumerate}

Las pruebas deberán registrar la versión del código, los artefactos utilizados,
los resultados y las incidencias. Se emplearán datos de prueba o una selección
autorizada, evitando alterar conjuntos de evaluación reservados. El éxito de
un recorrido funcional no sustituye la evaluación científica de la calidad
del clasificador ni la medición de los requisitos de rendimiento.

\subsection{Puesta en marcha local documentada}

Con corte al 4 de noviembre de 2026, el proyecto dispone de un nodo local
configurado con PostgreSQL en Cloud SQL, almacenamiento en GCS y archivos
persistentes propios. La actualización del 25 de octubre aplicó la revisión
SQL \texttt{0004\_integration}, conservó la historia previa y comprobó el
acceso, los permisos y el catálogo de modelos. También se restauraron copias
lógicas anterior y posterior en bases locales separadas, comparando tablas
y secuencias.

El catálogo verificado en ese mantenimiento incluyó un ensamble de aves como
predeterminado y cinco modelos individuales. La preparación no acreditó la
disponibilidad del ensamble industrial, cuyos pesos seguían pendientes. La
operación local de aves constituye un punto de partida para la implantación y
no permite extender sus resultados a otros dominios sin comprobaciones propias.

El 1 de noviembre se documentó la reanudación de los servicios y la
comprobación del acceso administrador. Ese registro no estableció la causa de
su detención previa ni verificó un mecanismo permanente de arranque. El 4 de
noviembre se comprobó la cuenta y el permiso técnico de un segundo integrante
del equipo; el arranque de su nodo Linux permaneció pendiente.

La Tabla~\ref{tab:estado-implantacion} resume el alcance de estos antecedentes.

\begin{table}[htbp]
    \centering
    \caption{Estado documentado de la implantación al 4 de noviembre de 2026.}
    \label{tab:estado-implantacion}
    \begingroup
    \small
    \renewcommand{\arraystretch}{1.18}
    \begin{tabular}{p{0.22\linewidth} p{0.68\linewidth}}
        \hline
        \textbf{Aspecto} & \textbf{Alcance observado o pendiente} \\
        \hline
        Nodo local & Configuración oficial, conexiones cloud, acceso y catálogo comprobados en el equipo del proyecto. \\
        Recuperación SQL & Respaldos lógicos restaurados y contrastados en bases locales separadas. Política externa y recuperación administrada pendientes. \\
        Segundo equipo & Cuenta web y permiso técnico comprobados; configuración y arranque Linux pendientes de verificación. \\
        Capacitación & Plan formativo propuesto; ejecución y evaluación con participantes pendientes. \\
        Entorno institucional & Servidor, acceso HTTPS, servicio permanente y política de continuidad pendientes de preparación y aceptación. \\
        \hline
    \end{tabular}
    \endgroup
\end{table}

\subsection{Continuidad operacional y transferencia institucional}

La transferencia a la institución deberá formalizar quién administra las
cuentas, quién resuelve incidentes y quién aplica actualizaciones. Se
prepararán el origen web autorizado, HTTPS y la supervisión de los procesos,
incluido su comportamiento después de reiniciar el equipo.

El respaldo deberá considerar la base de datos y los archivos operativos
necesarios para recuperar el nodo. Se definirán destino, periodicidad,
retención y responsable de las copias, y se ensayará la restauración en un
entorno separado. El respaldo de un checkpoint no reemplaza la recuperación
de los registros SQL ni de los índices de datos.

Antes de una actualización se conservarán una copia verificada de la base y
la identificación del software y de los artefactos previos. Una reversión
deberá utilizar código y esquema compatibles; cambiar únicamente la versión
del repositorio no revierte una migración SQL.

La aceptación institucional se realizará cuando el entorno de destino
complete el recorrido operacional, el personal designado demuestre las
competencias correspondientes y se disponga de procedimientos de soporte y
recuperación comprobados. La importación y exportación ZIP se habilitarán
únicamente donde estén disponibles Linux/Docker y los workers requeridos.
Los resultados de esta aceptación deberán incorporarse al registro de puesta
en marcha cuando se ejecute la transferencia.
